\section{The finite Artin matrix and the final contradiction}
\label{sec:assembly}

We now combine the length restriction of Proposition~\ref{prop:braid-angle}
with the obstruction of Proposition~\ref{prop:triangle-obstruction}.
The additional generators in the following matrix make the group $T$ a
common centralizer.  This will produce a geometric action of $T$ from
any hypothetical geometric action of the full group.

\begin{definition}\label{def:matrix}
Let $\mathcal E=\{(a,b),(b,c),(c,a)\}$, as before, and set
\[
 S=\{a,b,c,d,e\}\ \sqcup
   \{s_j^{uv}:(u,v)\in\mathcal E,\ 3\le j\le39\}.
\]
All symbols in the second set are distinct.  For each $(u,v)\in\mathcal E$,
also write $s_1^{uv}=u$ and $s_2^{uv}=v$; these are aliases for existing
elements of $S$.  Define a symmetric matrix $M=(m_{st})_{s,t\in S}$ by
the following prescriptions:
\begin{enumerate}
\item $m_{ss}=1$ for every $s\in S$, and
      $m_{ab}=m_{bc}=m_{ca}=3$.
\item In each block $\{s_1^{uv},\ldots,s_{39}^{uv}\}$, set
      \[
      m_{s_i^{uv},s_j^{uv}}=
      \begin{cases}
      3,&|i-j|=1,\\
      2,&|i-j|\ge2.
      \end{cases}
      \]
\item For $q\in\{d,e\}$ and $u\in\{a,b,c\}$, set $m_{qu}=2$.
\item Every remaining off-diagonal entry is $\infty$.
\end{enumerate}
Let $G=A_M$ be the resulting Artin group, with the generators denoted
by the elements of $S$ themselves.  Let $G_0$ be the group presented
by the generators $S\setminus\{d,e\}$ and all their prescribed relations.
\end{definition}

Two distinct blocks meet in just one generator, so their prescriptions
do not conflict.  Each pair among $a,b,c$ is the initial pair of exactly
one block, with label $3$ in both prescriptions.  Thus the definition
specifies exactly one matrix, on
\[
 |S|=5+3\cdot37=116
\]
generators.  In particular $m_{de}=\infty$, and neither $d$ nor $e$
has a finite label with any of the private block generators.

\subsection{The required embeddings}

\begin{lemma}\label{lem:pair-embeddings}
For every $(u,v)\in\mathcal E$, the homomorphism
\[
 B_3=\langle t_1,t_2\mid t_1t_2t_1=t_2t_1t_2\rangle
 \longrightarrow T,
 \qquad t_1\longmapsto u,\quad t_2\longmapsto v,
\]
is injective.  For every $l\ge3$, the homomorphism $B_3\to B_l$
given by $t_1\mapsto s_1$ and $t_2\mapsto s_2$ is also injective.
\end{lemma}

\begin{proof}
For the first assertion, let $w$ be the remaining element of
$\{a,b,c\}$.  The assignment
\[
 u\longmapsto t_1,\qquad v\longmapsto t_2,
 \qquad w\longmapsto t_1
\]
defines a homomorphism $T\to B_3$: the braid relation for $u,v$ is
the displayed defining relation of $B_3$, the relation for $v,w$ is
that relation with its sides interchanged, and the relation for $w,u$
becomes $t_1^3=t_1^3$.  Its composite with the asserted inclusion is
the identity of $B_3$.

For the second assertion we construct a word invariant, retaining
the entire permutation of the strands as part of its state.  Place
$l$ distinct tokens, named $1,\ldots,l$, in $l$ ordered positions.
Tokens $1,2,3$ are called selected.  A state consists of an ordered
list of all $l$ tokens and an element $q\in B_3$.  Read a signed
letter $s_i^{\varepsilon}$, $\varepsilon\in\{1,-1\}$, as follows.
Swap the tokens in positions $i,i+1$, irrespective of the sign.
If both are selected, they are consecutive in the ordered list of
selected tokens; if they occupy positions $j,j+1$ in that shorter
list, replace $q$ by $qt_j^{\varepsilon}$.  If at most one is selected,
leave $q$ unchanged.  The integer $j$ is computed before the swap,
and is either $1$ or $2$.

We verify the defining relations as transformations of all states,
not merely of the initial state.  The operation for $s_i^{-1}$ is
the inverse of that for $s_i$: a second swap restores the full token
list, and, when a letter was recorded, its selected-list index is
still $j$, so that the two recorded letters cancel.  For
$|i-k|\ge2$, the swaps at $i$ and $k$ involve disjoint pairs of
positions and have the same final token list in either order.
With just three selected tokens, at most one of these pairs consists
of two selected tokens.  Swapping the other pair does not change
the selected-list index of the recorded pair, because moving a
selected token past an unselected token preserves the order of the
selected list.  Hence the positive commutation relation is respected.

For a positive braid relation, name the three tokens initially in
positions $i,i+1,i+2$ by $A,B,C$.  The swaps in $s_is_{i+1}s_i$
exchange, in order, the pairs
\[
 (A,B),\ (A,C),\ (B,C),
\]
whereas those in $s_{i+1}s_is_{i+1}$ exchange them in the opposite
order.  Both finish with the ordered triple $C,B,A$ and fix all
other positions.  If at most one of $A,B,C$ is selected, neither
word records a letter.  If exactly two are selected, each word
records their single crossing with positive sign and the same
selected-list index: the third selected token is outside these
three consecutive positions and stays on the same side throughout.
If all three are selected, the two records are respectively
$t_1t_2t_1$ and $t_2t_1t_2$, which represent the same element of $B_3$.

Thus the invertible state transformations satisfy every defining
positive relation of $B_l$.  They consequently respect arbitrary
signed word equalities: each defining relator acts trivially, as
does its inverse or any conjugate, and free cancellations act
trivially by the inverse check.  This is precisely the normal-closure
description of equality in the presented group.  In particular,
the verification applies after every signed prefix, and the equal
full token lists ensure identical processing of every suffix.

Start with the token list $(1,\ldots,l)$ and $q=1$.  A word using
only $s_1^{\pm1},s_2^{\pm1}$ always swaps two of the first three
tokens.  Its record is exactly the same word with $s_i$ replaced
by $t_i$, including every sign.  If the original word is trivial
in $B_l$, its state transformation fixes this initial state, and
its record is therefore trivial in $B_3$.  This proves injectivity.
The construction asserts a state invariant; it does not assert that
discarding the other tokens gives a homomorphism $B_l\to B_3$.
\end{proof}

\subsection{Normal forms and the common centralizer}

We use the standard normal-form theorem for free products with
amalgamation; see Serre \cite[Chapter~I, \S1.2, Theorems~1--2]{Serre1980}.
The elementary proof below records the exact form needed here, including
the case in which multiplication removes the final representative.

\begin{lemma}\label{lem:amalgam-normal-form}
Let $H_1,H_2$ be groups with a common embedded subgroup $E$, and
let $H_1*_E H_2$ be the group obtained by identifying the two copies
of $E$ in their free product.  Each factor embeds in this group,
and their intersection is $E$.  A product $a_1\cdots a_n$, $n\ge1$,
whose factors belong alternately to $H_1\setminus E$ and
$H_2\setminus E$ is nontrivial.  If such a product contains a
factor from $H_i\setminus E$, it does not belong to the other factor.
\end{lemma}

\begin{proof}
For $i=1,2$, choose a set $R_i\subset H_i$ such that every element
of $H_i$ is uniquely $re$ with $r\in R_i$, $e\in E$, and with $1$
representing $E$.  Consider formal states
\[
 (r_1,\ldots,r_n;e),\qquad n\ge0,
\]
where every $r_j$ is a nonidentity representative and their factor
types alternate; for $n=0$ only the terminal element $e$ is present.
Representatives retain their factor types as part of the state.

Define right multiplication of a state by $h\in H_i$ as follows.
If its final representative has type $i$, isolate the tail $re$
with that representative $r$; otherwise use the missing
representative $r=1$.  The preceding list $P$ is empty or ends in
the other factor.  Decompose the element $reh$ of $H_i$ uniquely as
$r'e'$.  Replace the tail by $r';e'$, omitting $r'$ when $r'=1$.

For fixed $i$, these operations are a right action of $H_i$.
Indeed, after multiplying by $h\in H_i$, the prefix $P$ stays fixed.
For a further $k\in H_i$, the tail is rewritten from
$r'e'k=(reh)k$.  If $r'=1$ was omitted, the rule simply uses the
missing representative $1$ and rewrites $e'k$, with exactly the
same result.  Uniqueness of the decomposition in $H_i$ and
associativity therefore give the operation for $hk$.  Multiplication
by $1$ fixes every state, and multiplication by $h^{-1}$ is inverse
to multiplication by $h$.

If $h\in E$, either factor action leaves all representatives fixed
and changes only the terminal element from $e$ to $eh$.  The two
actions therefore agree on $E$ and give a right action of the
presented amalgam: all multiplication relations within the factors
and all identifications of $E$ act identically.

Every state is reached from $(\,;1)$ by reading the word
$r_1\cdots r_ne$.  Conversely, applying the algorithm to an arbitrary
word uses only valid equalities in the amalgam and produces a word
of that form.  If two state words were equal in the amalgam, their
actions on $(\,;1)$ would give the same state.  This proves both
existence and uniqueness of these normal forms.  In particular,
the normal form of an element of $H_i$ is its single representative
and terminal $E$-element, or just its $E$-element; this proves the
factor embeddings and the intersection assertion.

Finally, when processing an alternating product $a_1\cdots a_n$
with every $a_j\notin E$, at each step the terminal $E$-element
times $a_j$ is still outside $E$.  Since its factor type differs
from that of the preceding representative, one new nonidentity
representative is appended.  Its normal form thus has exactly $n$
representatives, of the indicated alternating types.  It is
nontrivial and cannot be the normal form of an element of a factor
whose type omits any of these representatives.
\end{proof}

Write $B_{40}^{uv}$ for the abstract braid group in the block
corresponding to $(u,v)$, and $B_3^{uv}$ for its initial two-generator
subgroup.  By Lemma~\ref{lem:pair-embeddings}, each $B_3^{uv}$ embeds
both in that block and in $T$.  Applying
Lemma~\ref{lem:amalgam-normal-form} successively therefore gives
\begin{equation}\label{eq:iterated-amalgam}
 G_0\cong
 \bigl((T*_{B_3^{ab}}B_{40}^{ab})
       *_{B_3^{bc}}B_{40}^{bc}\bigr)
       *_{B_3^{ca}}B_{40}^{ca}.
\end{equation}
At every stage $T$ remains embedded, so the next attaching subgroup
is embedded as well.  The presentations agree: identifying the
first two generators of each block with its specified pair in $T$
leaves exactly the generators and relations defining $G_0$.
Thus $T$ and all three braid blocks are embedded in $G_0$.

Let $F(d,e)$ be the free group on $d,e$, realized as reduced words
in $d^{\pm1},e^{\pm1}$ with multiplication by concatenation and
free cancellation, and put
$B=T\times F(d,e)$.  The direct product has the presentation for $T$
together with the two free generators and the six commutation
relations between $\{d,e\}$ and $\{a,b,c\}$.  To check this description
directly, the commutations move every $T$-letter to the left of every
free letter; the evident homomorphism to the direct product and
the two commuting factor homomorphisms are inverse on the generators.
Comparison with Definition~\ref{def:matrix} now gives
\begin{equation}\label{eq:full-amalgam}
 G\cong G_0*_T B
       =G_0*_T\bigl(T\times F(d,e)\bigr).
\end{equation}
Lemma~\ref{lem:amalgam-normal-form} makes both factors, and hence
$T$ and the three braid blocks, embedded subgroups of $G$.

\begin{lemma}\label{lem:common-centralizer}
For the elements $d,e$ of $G$,
\[
 C_G(d)\cap C_G(e)=T.
\]
\end{lemma}

\begin{proof}
Every element of $T$ commutes with $d$ and $e$ by the presentation.
Conversely, first suppose that $g\notin B$.  Its normal form in
\eqref{eq:full-amalgam} can be written $g=wb$, where $b\in B$ and
$w$ is a nonempty alternating product of elements outside $T$,
whose last factor lies in $G_0\setminus T$.  Indeed, absorb the
terminal $T$-element and, if present, the last $B$-representative
into $b$.  There is at least one $G_0$-representative because
$g\notin B$.

Under the projection $B\to F(d,e)$, the element $bdb^{-1}$ maps
to a conjugate of the nonidentity free generator $d$.  Consequently
$bdb^{-1}\in B\setminus T$.  The word
\[
 gdg^{-1}=w(bdb^{-1})w^{-1}
\]
is reduced in the sense of Lemma~\ref{lem:amalgam-normal-form}:
its middle factor lies in $B\setminus T$, both neighbors lie in
$G_0\setminus T$, and all remaining adjacent factors alternate.
It contains a $G_0\setminus T$ factor and hence cannot belong to
$B$.  In particular it cannot equal $d$.  Thus $C_G(d)\subset B$.

It remains to compute the common centralizer within $B$.
If a reduced word $f\in F(d,e)$ commutes with $d$, write
$f=v d^q$, where $q\in\mathbb Z$ and either $v$ is empty or its
last letter is $e$ or $e^{-1}$.  This is obtained by removing the
maximal terminal string of $d$-letters from the reduced word.
Then
\[
 fdf^{-1}=v d v^{-1}.
\]
If $v$ is nonempty, this displayed word is freely reduced, has
length $2|v|+1>1$, and cannot equal $d$.  Hence $f=d^q$.
For $q\ne0$, the word $d^q e d^{-q}$ is freely reduced and is
not $e$, so only the identity commutes with both free generators.
An element $(t,f)\in T\times F(d,e)$ therefore centralizes both
$d$ and $e$ exactly when $f=1$.  This proves the assertion.
\end{proof}

\subsection{A geometric action of the common centralizer}

Geometric actions of centralizers on minimum-displacement sets are
classical.  Swenson \cite[Theorem~13]{Swenson1999} credits this result
to Ruane \cite{Ruane2001}; Swenson also studies intersections of convex
subgroups, and his proof of Theorem~17 treats simultaneous centralizers
\cite[Theorems~16--17]{Swenson1999}.  The following finite-tuple
compactness argument gives the precise common-displacement-sublevel
formulation needed here.  It applies in every dimension and does not
require the elements being centralized to commute.

\begin{lemma}\label{lem:geometric-centralizer}
Suppose a group $H$ acts geometrically on a nonempty proper CAT(0)
space $X$.  For finitely many elements $q_1,\ldots,q_m\in H$,
$m\ge1$, set
\[
 C=\bigcap_{j=1}^m C_H(q_j).
\]
If $A\ge0$ is finite and the set
\[
 Y_A=\{x\in X:d_X(x,q_jx)\le A\text{ for }1\le j\le m\}
\]
is nonempty, then $Y_A$, with its inherited metric, is a proper
CAT(0) space, invariant under $C$, and the restricted $C$-action
is geometric.  Every element of $C$ has the same stable translation
length in $Y_A$ as in $X$.
\end{lemma}

\begin{proof}
Each function $x\mapsto d_X(x,q_jx)$ is continuous and convex.
For the convexity, if $\gamma$ is the geodesic from $x_0$ to $x_1$
parametrized on $[0,1]$, then $q_j\gamma$ is the geodesic from
$q_jx_0$ to $q_jx_1$, and Lemma~\ref{lem:comparison} gives
\[
 d_X(\gamma(t),q_j\gamma(t))
 \le (1-t)d_X(x_0,q_jx_0)+t d_X(x_1,q_jx_1).
\]
Thus $Y_A$ is closed and convex.  Its geodesic triangles and their
metrics are inherited from $X$, so it is CAT(0); each closed bounded
subset of $Y_A$ is closed and bounded in $X$ and is therefore compact.
If $c\in C$, then
\[
 d_X(cx,q_jcx)=d_X(cx,cq_jx)=d_X(x,q_jx),
\]
so $Y_A$ is $C$-invariant.  Compact subsets of $Y_A$ are compact
in $X$, hence the restricted action satisfies the required
compact-set properness condition.

It remains to prove cocompactness.  Choose compact $K\subset X$
with $HK=X$, and let
\[
 L=\{z\in X:d_X(z,K)\le A\}.
\]
This closed bounded neighborhood of $K$ is compact, since $X$ is
proper.  Accordingly
\[
 \mathcal F=\{r\in H:rL\cap L\ne\varnothing\}
\]
is finite.  For any $h\in H$ such that $hK\cap Y_A\ne\varnothing$,
choose $p\in K$ with $hp\in Y_A$.  For every $j$,
\[
 d_X(p,h^{-1}q_jhp)=d_X(hp,q_jhp)\le A.
\]
Both $p$ and $h^{-1}q_jhp$ belong to $L$, so
$h^{-1}q_jh\in\mathcal F$.  Only finitely many ordered tuples
\[
 (h^{-1}q_1h,\ldots,h^{-1}q_mh)
\]
are therefore possible.  Choose one representative $h_i$ for each
tuple that occurs, with $1\le i\le N$.

If $h$ has the same tuple as $h_i$, then, for every $j$,
\[
 h^{-1}q_jh=h_i^{-1}q_jh_i
 \quad\Longrightarrow\quad
 (hh_i^{-1})q_j(hh_i^{-1})^{-1}=q_j.
\]
Consequently $hh_i^{-1}\in C$, or equivalently
\[
 h\in C h_i,\qquad h=ch_i\quad\text{for some }c\in C.
\]
Here $C h_i$ is a \emph{right} coset, in the standard convention.
Now set
\[
 K_C=Y_A\cap\bigcup_{i=1}^N h_iK.
\]
This set is compact.  Given $x\in Y_A$, write $x=hp$ with $p\in K$,
choose its tuple representative $h_i$, and write $h=ch_i$ as above.
Then $c^{-1}x=h_ip$ belongs to both $Y_A$ and $h_iK$, so $x\in CK_C$.
Thus $CK_C=Y_A$, proving cocompactness.

Finally, for $y\in Y_A$, $c\in C$, and every positive integer $n$,
invariance and the inherited metric give
\[
 d_{Y_A}(y,c^ny)=d_X(y,c^ny).
\]
Dividing by $n$ and taking limits proves equality of stable lengths.
\end{proof}

\begin{proof}[Proof of Theorem~\ref{thm:main}]
Let $M$ be the explicit finite matrix of Definition~\ref{def:matrix}.
Suppose, for a contradiction, that $G=A_M$ acts properly and
cocompactly by isometries on a nonempty proper CAT(0) space $X$.
No restriction is placed on the dimension of $X$.

Choose $x_0\in X$ and a finite number
\[
 A\ge\max\{d_X(x_0,dx_0),d_X(x_0,ex_0)\}.
\]
The set
\[
 Y=\{x\in X:d_X(x,dx)\le A,\ d_X(x,ex)\le A\}
\]
contains $x_0$.  By Lemmas~\ref{lem:common-centralizer} and
\ref{lem:geometric-centralizer}, $T$ acts geometrically on $Y$,
and its stable lengths there agree with its stable lengths in $X$.

Sending every Artin generator of $G$ to $1\in\mathbb Z$ defines
a homomorphism $G\to\mathbb Z$: both sides of each defining
Artin relation have the same number of letters.  In particular,
every squared generator maps to $2$ and has infinite order.
Lemma~\ref{lem:semisimple} therefore gives
$\ell_X((s_1^{uv})^2)>0$ for each $(u,v)\in\mathcal E$.
The corresponding embedded $B_{40}$ block satisfies all the
relations required by Proposition~\ref{prop:braid-angle}, so that
\[
 \frac{\ell_X((uv)^3)}{3\ell_X(u^2)}
 =\frac{\ell_X((uv)^3)}{3\ell_X(v^2)}
 >\frac12
 \qquad ((u,v)\in\mathcal E).
\]
Equality of stable lengths transfers all three strict inequalities
to the geometric $T$-action on $Y$.  This contradicts
Proposition~\ref{prop:triangle-obstruction}.

The $116$-generator Artin group $A_M$ therefore has no proper
cocompact isometric action on any nonempty proper CAT(0) space.
This is the required counterexample.
\end{proof}
